\documentclass[10pt]{article}

% Specify the margins and text size.
\setlength{\textwidth}{6.8in}
\setlength{\textheight}{9.5in}
\setlength{\oddsidemargin}{0pt}
\setlength{\evensidemargin}{0pt}
\setlength{\topmargin}{0pt}
\setlength{\hoffset}{-0.15in}
\setlength{\voffset}{-1in}

\setlength{\parskip}{5pt}
\setlength{\parindent}{0pt}

% Load some fonts and symbol packages
\usepackage{latexsym}
\usepackage{pifont}       % contains 'star' symbol for counterinsurgency handbook title
\usepackage{yfonts} 
\usepackage{amsmath}
\usepackage{amsfonts}

\usepackage{graphicx}     % actually, this is loaded by pstricks
\usepackage[T1]{fontenc}
\usepackage{ifthen}
\usepackage{pstricks,pst-grad,pst-text,pst-node,multido,pst-plot,calc,pst-3dplot}
%\usepackage[all]{xy}
%\usepackage{animate}

% The hyperref package inserts links into the pdf.
\definecolor{MyLinkColor}{rgb}{.1,.2,1}
\definecolor{MyCiteColor}{rgb}{.1,1,.2}
\definecolor{MyURLColor}{rgb}{.4,.4,.4}
\usepackage[backref=true,pagebackref=false,hyperindex,colorlinks=true,
  linkcolor=MyLinkColor,urlcolor=MyURLColor]{hyperref}


% The tweaklist package is something I found on the web.  It provides a simple interface
% for making changes to spacing used in the itemize and enumerate environments.  Comment
% this out if you don't care to use tweaklist.
\usepackage{tweaklist}
\renewcommand{\itemhook}{\setlength{\parskip}{2pt}\setlength{\parsep}%
{1pt}\setlength{\topsep}{0pt}\setlength{\itemsep}{0pt}}

\newcommand{\U}{\underline{\hspace{5pt}}}

\usepackage{listings}
\newcommand{\Z}{\hphantom{0}}


\newcommand{\zobox}[2]{\begin{pspicture}(-0.2,0.2)(2.4,1.3)
\rput[r](-0.08,0.35){\small #1\vphantom{,'}}\psframe(0.0,0.05)(0.5,0.65)\rput(0.25,0.35){0}
   \rput[l](0.58,0.35){s\vphantom{',1}}
\rput[r](1.74,0.35){\small #2\vphantom{,'}}\psframe(1.8,0.05)(2.3,0.65)\rput(2.05,0.35){1}
   \rput[l](2.38,0.35){s\vphantom{',1}}
%\psframe(1,0.05)(1.6,0.65)\rput(1.3,0.35){1}\rput[bl](1.64,0.64){\scriptsize #2}
\psline(-0.9,0.9)(-0.9,-0.2)(2.8,-0.2)(2.8,0.9)(-0.9,0.9)
\end{pspicture}}

\newcommand{\threebox}[3]{\begin{pspicture}(0.2,0.2)(2.5,1.1)
\psframe(0,0)(0.7,0.7)\rput(0.35,0.35){#1}
\psframe(1,0)(1.7,0.7)\rput(1.35,0.35){#2}
\psframe(2,0)(2.7,0.7)\rput(2.35,0.35){#3}
\psline(-0.3,0.9)(-0.3,-0.2)(3,-0.2)(3,0.9)(-0.3,0.9)
\end{pspicture}}

\begin{document}
\pagestyle{empty}
\lstset{language=R, showspaces=false, showstringspaces=false}

\href{http://www.su.edu}{\includegraphics[height=1.75cm]{sulogo.eps}}
\vspace{-1.79cm}

{{\ }\hfill\small
\begin{tabular}{cl}
& Math 207\\
& Introduction to Statistics\\
\end{tabular}
}
\setlength{\baselineskip}{1.05\baselineskip}

\begin{center}
\textbf{\large Expected Values and Standard Errors}
\end{center}

The \textbf{mean} of a list of numbers $x_1$, $\dots$, $x_n$ is\vspace{-4pt}
\[\mbox{mean} = \frac{1}{n}\,(x_1 + x_2 + \cdots + x_n) = \frac{1}{n}\,\sum_{i=1}^n\,x_i.\vspace{-4pt}\]
The \textbf{population standard deviation} is\vspace{-4pt}
\[\mbox{SD} = \sqrt{\frac{1}{n}\,\big((x_1-\mbox{mean})^2 + \cdots + (x_n - \mbox{mean})^2\big)}\vspace{-4pt}\]
The \textbf{sample standard deviation} is\vspace{-4pt}
\[\mbox{sd} = \sqrt{\frac{1}{n-1}\,\big((x_1-\mbox{mean})^2 + \cdots + (x_n - \mbox{mean})^2\big)}\vspace{-4pt}\]

1. When we are working with a box model and know what all the tickets are, then we should use
the population standard deviation. Compute the mean and SD for each box.

\hspace{10pt} a) Flipping a Coin:\hspace{.25in}
\begin{pspicture}(0.2,0.2)(2.5,1.1)
\psframe(0,0)(0.7,0.7)\rput(0.35,0.35){0}
\psframe(1,0)(1.7,0.7)\rput(1.35,0.35){1}
\psline(-0.3,0.9)(-0.3,-0.2)(2,-0.2)(2,0.9)(-0.3,0.9)
\end{pspicture}
\vspace{1.3in}

\hspace{10pt} b) Flipping an Unfair  Coin:\hspace{.25in}
\begin{pspicture}(0.2,0.2)(2.5,1.1)
\psframe(0,0)(0.7,0.7)\rput(0.35,0.35){0}
\psframe(1,0)(1.7,0.7)\rput(1.35,0.35){1}
\psframe(2,0)(2.7,0.7)\rput(2.35,0.35){1}
\psline(-0.3,0.9)(-0.3,-0.2)(3,-0.2)(3,0.9)(-0.3,0.9)
\end{pspicture}
\vspace{1.5in}


\hspace{10pt} c) Rolling a die:\hspace{.4in}
\begin{pspicture}(0.2,0.2)(2.5,1.1)
\psframe(0,0)(0.7,0.7)\rput(0.35,0.35){1}
\psframe(1,0)(1.7,0.7)\rput(1.35,0.35){2}
\psframe(2,0)(2.7,0.7)\rput(2.35,0.35){3}
\psframe(3,0)(3.7,0.7)\rput(3.35,0.35){4}
\psframe(4,0)(4.7,0.7)\rput(4.35,0.35){5}
\psframe(5,0)(5.7,0.7)\rput(5.35,0.35){6}
\psline(-0.3,0.9)(-0.3,-0.2)(6,-0.2)(6,0.9)(-0.3,0.9)
\end{pspicture}
\vfill
\eject

\hspace{10pt} d) A Box with Two Different Numbers:\hspace{.25in}
\begin{pspicture}(0.2,0.2)(2.5,1.1)
\psframe(0,0)(0.7,0.7)\rput(0.35,0.35){1}
\psframe(1,0)(1.7,0.7)\rput(1.35,0.35){9}
\psline(-0.3,0.9)(-0.3,-0.2)(2,-0.2)(2,0.9)(-0.3,0.9)
\end{pspicture}
\vspace{1.5in}


\hspace{10pt} e) A Zero-One Box:\hspace{.25in} \zobox{530}{470}
\vspace{1.6in}


2. Write down a box model for the following random experiments.

\hspace{10pt} a) Flipping a fair coin ($0=\mbox{tails}$ and $1=\mbox{heads})$
\vspace{.7in}

\hspace{10pt} b) Flipping a coin for which the chance of heads is $2/3$.
\vspace{.8in}

\hspace{10pt} c) Rolling a fair die.
\vspace{1in}


\hspace{10pt} d) Rolling a fair die and counting the number of times you get a 5.
\vspace{1in}

\hspace{10pt} e) Rolling a fair die and counting the number of times you get a 3 or a 5.
\vfill
\eject


\end{document}
